\documentclass[12pt]{article}
\usepackage[utf8]{inputenc}
\usepackage[T1]{fontenc}
\usepackage{hyperref}
\title{Tay Language Resources for K-12 Curriculum Development}
\author{Compiled by Claude Sonnet 5 (claude-sonnet-5) for Andr\'as Kornai}
\date{}
\begin{document}
\maketitle
\section{Introduction}
Tay (also called Tai; ISO 639-3 taw; Glottocode taii1241) is a Trans-New Guinea language of the Madang branch, Kalam group, spoken in Madang Province, Papua New Guinea. Speaker estimates differ: about 900 (Wikipedia, 1990), about 1,000 (MPI numeral data) and about 2,800 (Joshua Project). It should not be confused with Tayap or with Thai-type Tai languages.
\section{Primary Educational Resources}
No primers, textbooks or school programmes were found in search results.
\section{Linguistic and Reference Materials}
\begin{itemize}
\item Wikipedia, Tai language (New Guinea): \url{https://en.wikipedia.org/wiki/Tai_language_(New_Guinea)}
\item Max Planck (EVA) numeral system page for Tay, with number words and tally system: \url{https://lingweb.eva.mpg.de/channumerals/Tay.htm}
\item Max Planck (EVA) numeral page, Tai (Tay): \url{https://lingweb.eva.mpg.de/channumerals/Tai-Tay-NPG.htm}
\end{itemize}
\section{Religious and Community Resources}
Joshua Project reports that a New Testament was published in 2013 and that Global Recordings Network audio teaching exists: \url{https://joshuaproject.net/people_groups/19111}. I did not verify the scripture itself.
\section{Digital and Audio Resources}
Audio Bible teaching is mentioned on the Joshua Project page above. No other digital resources were found.
\section{Curriculum Development Context}
Documentation is sparse. The New Testament, if accessible, could provide written text. The language is used by all community members as a first language, so oral materials would be the starting point. Community consultation on orthography is needed.
\end{document}